%!TEX root = TQFTnotes.tex
%%%%%%%%%%%%%%%%%%%%%%%%%%%%%%
\chapter*{Part 2 overview}\label{c:p2-overview}
This part is the heart of the book. Here we start the systematic study of TQFTs,
classifying and studying TQFTs in dimension 2. We also introduce the main tools
that will be used in all dimensions.

In \chref{c:2d-classification}, we discuss classification of 2-dimensional TQFT,
showing that such theories are classified by commutative Frobenius algebras. Our
approach is based on the use of Morse functions and families of Morse functions
(Cerf theory).

In \chref{c:frobenius}, we study Frobenius algebras, in particular classifying
all semisimple commutative Frobenius algebras and describing the corresponding
2d TQFTs.

In \chref{c:DW}, we describe the fundamental example of a TQFT: gauge theory
with finite gauge group, also known as Dijkgraaf--Witten theory. This theory can
be described in terms of $G$-covers of surfaces; we can also relate it to the
path integral description of TQFT.

After that, we move to describing extended TQFTs, i.e., theories where we
consider not just cobordisms but manifolds with corners. To give an accurate
definition of such an extended TQFT, one needs to introduce the language of
(weak) 2-categories, which we do in \chref{c:2categories}. We also introduce the
notion of a symmetric monoidal 2-category. Our key examples of 2-categories are
the 2-category of small categories (or a variation of it, the 2-category of all
locally finite $\kk$-linear abelian categories) and the category $\Alg$, whose
objects are algebras over $\kk$ and 1-morphisms are bimodules. Both
2-categories play an important role in future constructions.

Next, in \chref{c:duality-2cat} we discuss dualities in 2-categories, discussing
both the notion of a dual object (which only makes sense in a monoidal
2-category) and of a dual 1-morphism; in particular, we show that for the
2-category of small categories, where 1-morphisms are functors between
categories, so defined duality coincides with the classical notion of
adjointness for functors.
In \chref{c:separable-frobenius}, we study the duality in the 2-category $\Alg$
and show that fully dualizable objects in $\Alg$ are the separable algebras;
moreover, fixing a duality between $\ev$ and $\coev$ 1-morphisms is equivalent
to choosing a structure of a Frobenius algebra, which provides the link to the
construction of (non-extended) TQFTs done earlier.


After completing all these preliminaries, we finally define the notion of an
extended 2d TQFT in \chref{c:2d-extended}. Following work of Schommer--Pries, we
construct the set of generators and relations for the 2-category $\Bord_2$ of
2-dimensional bordisms and use it to show that extended 2d TQFTs are classified
by separable symmetric Frobenius algebras.

Finally, in \chref{c:PL}--\chref{c:2d-lattice} we present an alternative
approach to construction of extended 2d TQFTs. Instead of using Morse functions
to create a presentation of a cobordism as a composition of elementary
cobordisms, we use triangulations (or more generally, PL cell decompositions);
instead of Cerf moves, we use Pachner moves. We show how in dimension 2, it
gives a very explicit construction of an extended 2d TQFT from a separable
symmetric Frobenius algebra $A$, providing a different way of proving the same
classification result.